\chapter{Background}

\section{Causal graphs and conditional independence}

Let \(V\) be a finite set of variables. A graph over \(V\) is a pair \(G=(V,E)\), where \(E \subseteq \{(x,y) \in V \times V \mid x \neq y\}\) is a set of edges. A directed edge from \(x\) to \(y\) is written \(x \to y\). Two variables \(x,y \in V\) are adjacent in \(G\) if either \(x \to y\) or \(y \to x\) belongs to \(E\). A path, in the sense used here, is a sequence of distinct nodes \(x_1,\ldots,x_n\) such that each consecutive pair is adjacent. A path is directed if \(x_i \to x_{i+1}\) for all \(1 \leq i < n\). A directed cycle is a sequence \(x_1,\ldots,x_n,x_1\), with \(x_1,\ldots,x_n\) distinct, such that \(x_i \to x_{i+1}\) for \(1 \leq i < n\) and \(x_n \to x_1\). A directed acyclic graph (DAG) is a directed graph with no directed cycles. In a causal graph, a directed edge \(x \to y\) represents the hypothesis that \(x\) is a direct cause of \(y\), relative to the variables included in \(V\) \cite{russo2024argumentative,glymour2019review}.

A central relation in graphical causal models is d-separation. Let \(p\) be a path in a DAG \(G\). A non-endpoint node \(z\) on \(p\) is a collider on \(p\) if the two adjacent edges on the path both point into \(z\), giving a subpath of the form \(x \to z \leftarrow y\). A non-endpoint node that is not a collider is a non-collider. A node \(d\) is a descendant of \(z\) if there is a directed path from \(z\) to \(d\). Given a conditioning set \(Z \subseteq V \setminus \{x,y\}\), a path \(p\) is active given \(Z\) if every non-collider on \(p\) is outside \(Z\), and every collider on \(p\) is either in \(Z\) or has a descendant in \(Z\). If no path between \(x\) and \(y\) is active given \(Z\), then \(x\) and \(y\) are d-separated given \(Z\) in \(G\), written \(x \indep_G y \mid Z\). If at least one path between \(x\) and \(y\) is active given \(Z\), then \(x\) and \(y\) are d-connected given \(Z\) in \(G\), written \(x \nindep_G y \mid Z\).

The corresponding statistical notation is written without the graph subscript. For variables \(x,y \in V\) and a conditioning set \(Z \subseteq V \setminus \{x,y\}\), \(x \indep y \mid Z\) denotes that \(x\) and \(y\) are conditionally independent given \(Z\) in the probability distribution over \(V\). Similarly, \(x \nindep y \mid Z\) denotes conditional dependence. When \(Z=\emptyset\), unconditional independence is written as \(x \indep y\). More generally, for pairwise disjoint sets \(X,Y,Z \subseteq V\), the notation \(X \indep Y \mid Z\) denotes that the variables in \(X\) are conditionally independent of the variables in \(Y\) given \(Z\).

The link between a DAG and a probability distribution is given by assumptions that connect graphical separation with statistical independence. Under the Markov condition, d-separation in the graph implies conditional independence in the distribution, so \(x \indep_G y \mid Z \Longrightarrow x \indep y \mid Z\). Faithfulness gives the converse direction: every conditional independence in the distribution is assumed to arise from d-separation in the graph, so \(x \indep y \mid Z \Longrightarrow x \indep_G y \mid Z\). Together, these assumptions allow conditional-independence evidence to be interpreted as constraints on possible causal graphs. In the causal-discovery setting considered here, one also typically assumes causal sufficiency: all common causes of the observed variables are included in \(V\).

A DAG \(G\) is compatible with an independence statement \(x \indep y \mid Z\) if \(x\) and \(y\) are d-separated given \(Z\) in \(G\). Graphically, compatibility with a dependence statement \(x \nindep y \mid Z\) is taken to mean that \(x\) and \(y\) are d-connected given \(Z\) in \(G\); under the Markov and faithfulness assumptions, this corresponds to statistical dependence in the distribution. This notion of compatibility is used later when causal evidence is represented inside an argumentation framework. Conditional-independence information may still fail to identify a unique DAG: distinct DAGs can entail the same conditional independences and are then Markov equivalent. Equivalently, Markov-equivalent DAGs have the same skeleton and the same v-structures.

\section{Assumption-Based Argumentation}

ABA is a structured argumentation formalism in which arguments and attacks are generated from rules, assumptions and contraries. The presentation here follows the notation used in the ABA and ABA Learning literature, and focuses on the definitions needed later for Causal ABA and ABA Learning \cite{bondarenko1997abstract,toni2014tutorial}.

\subsection{ABA frameworks}

An ABA framework is a tuple \(\mathcal{F} = \langle \mathcal{L}, \mathcal{R}, \mathcal{A}, \overline{\cdot} \rangle\), where \(\mathcal{L}\) is a formal language, \(\mathcal{R}\) is a set of inference rules, \(\mathcal{A} \subseteq \mathcal{L}\) is a set of assumptions, and \(\overline{\cdot}\) is a contrary mapping assigning to each assumption \(\alpha \in \mathcal{A}\) a sentence \(\overline{\alpha} \in \mathcal{L}\). Assumptions represent defeasible premises: they may be used in arguments, but they can also be defeated if their contraries are derived.

A rule \(r \in \mathcal{R}\) has the form \(s_0 \leftarrow s_1,\ldots,s_m\), where \(s_0,s_1,\ldots,s_m \in \mathcal{L}\). The sentence \(s_0\) is the head of the rule, and \(s_1,\ldots,s_m\) are its body. When \(m=0\), the rule is written \(s_0 \leftarrow\) and is treated as a fact. Rules encode strict inferential structure, while assumptions encode defeasible information. The contrary mapping determines how assumptions can be challenged: deriving \(\overline{\alpha}\) gives a reason to reject the assumption \(\alpha\).

\subsection{Arguments, attacks and stable extensions}

An argument for a claim \(s \in \mathcal{L}\) is a derivation of \(s\) from a set of assumptions using the rules in \(\mathcal{R}\). If \(\Delta \subseteq \mathcal{A}\) is a set of assumptions and \(s\) is derivable from \(\Delta\) using rules in \(\mathcal{R}\), this is written \(\Delta \vdash_{\mathcal{R}} s\). Here, \(\Delta\) is the support of the argument and \(s\) is its conclusion.

Attacks are generated through contraries. An argument \(\Delta \vdash_{\mathcal{R}} s\) attacks an assumption \(\alpha \in \mathcal{A}\) when \(s = \overline{\alpha}\). Equivalently, a set of assumptions \(\Delta\) attacks an assumption \(\alpha\) if \(\Delta \vdash_{\mathcal{R}} \overline{\alpha}\). A set of assumptions \(\Delta\) attacks another set of assumptions \(\Gamma\) if it attacks some \(\gamma \in \Gamma\). A set \(\Delta \subseteq \mathcal{A}\) is conflict-free if it does not attack any of its own assumptions.

The semantics of an ABA framework specify which sets of assumptions can be accepted together. In the assumption-set presentation used here, a set \(\Delta \subseteq \mathcal{A}\) is closed if whenever \(\Delta \vdash_{\mathcal{R}} \alpha\) for some \(\alpha \in \mathcal{A}\), then \(\alpha \in \Delta\). A stable extension of \(\mathcal{F}\) is a closed, conflict-free set of assumptions that attacks every assumption outside it. In flat ABA frameworks closedness is automatic; it is stated explicitly here because later Causal ABA frameworks may be non-flat.

A claim \(s \in \mathcal{L}\) is accepted with respect to a stable extension \(\Delta\) if there exists some \(\Delta' \subseteq \Delta\) such that \(\Delta' \vdash_{\mathcal{R}} s\); write \(\mathcal{F} \models_{\Delta} s\) when this holds. Assuming at least one stable extension exists, a claim is cautiously accepted if it is accepted with respect to every stable extension. It is bravely accepted if it is accepted with respect to at least one stable extension. These two acceptance modes are used later when describing ABA Learning: cautious acceptance requires support in every stable extension of the framework, whereas brave acceptance requires support in at least one stable extension \cite{de2023aba,de2024learning}.

\section{Causal ABA}

Causal ABA instantiates the ABA framework from the previous section for causal-graph reasoning. The basic idea is to represent candidate graph statements as defeasible assumptions, and to use ABA rules to enforce graph-theoretic constraints such as edge consistency, acyclicity and d-separation. Stable extensions can then be interpreted as compatible causal graph hypotheses \cite{russo2024argumentative}.

\subsection{Encoding causal claims in ABA}

Let \(V\) be a finite set of variables. For distinct \(x,y \in V\), Causal ABA uses an assumption \(\mathsf{arr}_{xy}\) to represent the claim that the graph contains the directed edge \(x \to y\). For each unordered pair \(\{x,y\}\), it uses one assumption \(\mathsf{noe}_{xy}\), read symmetrically, to represent that there is no edge between \(x\) and \(y\). These assumptions allow the framework to reason defeasibly about the possible edge status of each pair of variables.

The Causal ABA rules enforce consistency among the possible edge states. For a pair of distinct variables \(x,y \in V\), the alternatives \(\mathsf{arr}_{xy}\), \(\mathsf{arr}_{yx}\), and \(\mathsf{noe}_{xy}\) are mutually incompatible: the graph cannot contain both directions of the same edge, and it cannot both contain and not contain an edge between the same variables. This is represented in ABA by rules that derive the contrary of an assumption when an incompatible edge-state assumption is accepted. Under stable semantics, each extension selects exactly one of these three alternatives for each unordered pair.

Acyclicity is also encoded by ABA rules. If accepted arrows form a directed cycle, the corresponding cycle assumptions collectively attack each arrow in that cycle. Stable extensions therefore exclude sets of arrows that would create directed cycles. In this way, the accepted arrow assumptions determine a directed acyclic graph rather than an arbitrary directed graph.

Causal ABA also represents conditional-independence claims. For distinct \(x,y \in V\) and a conditioning set \(Z \subseteq V \setminus \{x,y\}\), an assumption of the form \(x \indep y \mid Z\) represents the defeasible claim that \(x\) and \(y\) are conditionally independent given \(Z\). The rules of the framework connect these assumptions to the graph induced by the accepted arrows. If the accepted arrow assumptions entail an active path between \(x\) and \(y\) given \(Z\), then \(x\) and \(y\) are d-connected in the induced graph, and the contrary of the independence assumption can be derived. In the full construction, active paths are represented using collider-trees, which encode ordinary paths together with branches from collider nodes to conditioned descendants. Thus, independence assumptions are attacked when they conflict with the d-separation structure of the graph.

Dependence information can also be integrated, but it requires additional machinery. A dependence fact \(x \nindep y \mid Z\) requires the induced graph to contain at least one \(Z\)-active path between \(x\) and \(y\). In the Causal ABA construction, this is not obtained merely by rejecting the corresponding independence assumption. Instead, dependence facts require additional blocked-path assumptions and rules, ensuring that if dependence is imposed as a fact, at least one path remains active. At the level of this report, the important point is that Causal ABA internalises both graph structure and conditional-independence information within a single ABA framework: arrows, no-edges and independence claims are represented defeasibly, while graph-theoretic constraints are encoded by rules.

\subsection{Extensions as compatible causal graphs}

A stable extension of a Causal ABA framework determines a graph through its accepted arrow assumptions. Given a stable extension \(\Delta\), define \(G_{\Delta} = (V,E_{\Delta})\), where \(E_{\Delta} = \{(x,y) \in V \times V \mid \mathsf{arr}_{xy} \in \Delta\}\). The edge-consistency and acyclicity rules ensure that \(G_{\Delta}\) is a DAG. Thus, stable extensions can be read as candidate causal graphs over the variable set \(V\).

The key correspondence is between accepted independence assumptions and d-separation in the graph induced by the extension. For stable extensions, the Causal ABA construction proves that \((x \indep y \mid Z) \in \Delta\) if and only if \(x \indep_{G_{\Delta}} y \mid Z\). That is, an independence claim is accepted exactly when it is supported by d-separation in the DAG determined by the same extension. Conversely, if the accepted arrows make \(x\) and \(y\) d-connected given \(Z\), then the corresponding independence assumption is defeated.

There is a distinction between assumptions available for selection by extensions and facts imposed from external evidence. Assumptions such as \(\mathsf{arr}_{xy}\), \(\mathsf{noe}_{xy}\), and \(x \indep y \mid Z\) are defeasible elements of the base Causal ABA framework; external causal-discovery information can be added as facts or constraints. When an independence or arrow claim is added as a fact, closure forces it to be included in every relevant extension. In ABA-PC, facts can also be handled with weights or weak constraints when test results are inconsistent, because erroneous statistical tests can otherwise yield no stable extensions. For example, a statistical procedure or domain expert may provide an independence claim, a dependence claim, an arrow claim, or a no-edge claim. When such information is imposed as fixed input, the relevant stable extensions are those whose induced graphs satisfy the added causal claims. The corresponding graph \(G_{\Delta}\) is then compatible with the external causal evidence when it satisfies the accepted independence, dependence and edge statements.

\section{ABA Learning}

ABA Learning studies how an ABA framework can be extended so that it accepts desired claims and does not accept undesired claims as consequences under the chosen acceptance mode. In this report, the relevant setting is learning from background knowledge together with positive and negative examples, under stable-extension semantics as used in the ASP-based ABA Learning and Causal ABA settings considered later. The section fixes the notation for this learning problem and introduces the transformation rules used to construct learnt ABA frameworks \cite{proietti2022learning,de2023aba}.

\subsection{Learning from examples}

Let \(\mathcal{F} = \langle \mathcal{L}, \mathcal{R}, \mathcal{A}, \overline{\cdot} \rangle\) be an initial ABA framework. This framework represents the background knowledge available before learning. Let \(E^+ \subseteq \mathcal{L}\) be a set of positive examples and \(E^- \subseteq \mathcal{L}\) be a set of negative examples, with \(E^+ \cap E^- = \emptyset\). A positive example is a claim that the learnt framework should accept. A negative example is a claim that the learnt framework should not accept.

An ABA Learning task is to construct an expanded ABA framework \(\mathcal{F}' = \langle \mathcal{L}', \mathcal{R}', \mathcal{A}', \overline{\cdot}^{\,\prime} \rangle\) such that the original framework is preserved and extended. Typically, \(\mathcal{L} \subseteq \mathcal{L}'\), \(\mathcal{R} \subseteq \mathcal{R}'\), and \(\mathcal{A} \subseteq \mathcal{A}'\). For existing assumptions \(\alpha \in \mathcal{A}\), the contrary mapping is preserved, so \(\overline{\alpha}^{\,\prime}=\overline{\alpha}\). The learnt framework may contain new rules, new assumptions, and contrary mappings for the new assumptions. Its purpose is to explain the examples under a chosen argumentation semantics. Some formulations also include a set \(T\) of learnable predicates, restricting which predicates may occur in learnt rule heads; this is omitted here when it is not needed.

Write \(\mathcal{F}' \models_{\mathrm{caut}} e\) when the claim \(e\) is cautiously accepted in \(\mathcal{F}'\), meaning that \(e\) is accepted with respect to every stable extension of \(\mathcal{F}'\). Write \(\mathcal{F}' \models_{\mathrm{brave}} e\) when \(e\) is bravely accepted in \(\mathcal{F}'\), meaning that \(e\) is accepted with respect to at least one stable extension of \(\mathcal{F}'\). These two acceptance modes give two corresponding forms of ABA Learning \cite{de2023aba,de2024learning}.

In cautious ABA Learning, a learnt framework \(\mathcal{F}'\) solves the task if \(\forall e \in E^+,\ \mathcal{F}' \models_{\mathrm{caut}} e\), and if \(\forall e \in E^-,\ \mathcal{F}' \not\models_{\mathrm{caut}} e\). That is, every positive example must be accepted with respect to every stable extension, and no negative example may be accepted with respect to every stable extension.

In brave ABA Learning, the requirement is existential over stable extensions: there must exist a stable extension \(\Delta\) of \(\mathcal{F}'\) such that \(\forall e \in E^+,\ \mathcal{F}' \models_{\Delta} e\), and \(\forall e \in E^-,\ \mathcal{F}' \not\models_{\Delta} e\). A negative example may still be accepted with respect to another stable extension without invalidating a brave solution. Thus, the learning objective is not only to derive the positive examples, but also to avoid accepting the negative examples under the selected consequence mode.

\subsection{Transformation rules}

ABA Learning constructs \(\mathcal{F}'\) by applying transformation rules to the initial framework. These rules add, generalise or revise rules and assumptions while preserving the background framework as the starting point. The main transformation rules used in this report are Rote Learning, Folding, Assumption Introduction and Fact Subsumption \cite{proietti2022learning,de2023aba,de2024learning}. The original ABA Learning framework also includes Equality Removal; it is omitted here because the later ASP/brave-learning algorithms used in this report focus on Rote Learning, Folding, Assumption Introduction and Fact Subsumption.

Rote Learning introduces example-specific rules. If \(p(t)\) is a positive example, Rote Learning may add a rule such as \(p(X) \leftarrow X=t\). This rule directly supports the example \(p(t)\). Rote Learning is useful because it can ensure that positive examples become derivable. By itself, however, it tends to memorise examples rather than express general patterns.

Folding generalises rules using existing background predicates. Suppressing the usual variable and equality side conditions, suppose the framework contains a rule of the form \(H \leftarrow B_1, B_2\), and the background knowledge contains a rule \(K \leftarrow B_1\). Folding can replace the occurrence of \(B_1\) in the first rule with \(K\), yielding a more abstract rule of the form \(H \leftarrow K, B_2\). The purpose of Folding is to move from example-specific rules towards intensional rules that reuse concepts already present in the background knowledge.

Assumption Introduction makes rules defeasible. A strict rule of the form \(H \leftarrow B\) may be transformed into \(H \leftarrow B, \alpha\), where \(\alpha\) is a new assumption. The transformed rule now applies only when \(\alpha\) is accepted. Exceptions can then be represented by rules that derive the contrary \(\overline{\alpha}\). This represents exceptions as undercutting attacks: deriving \(\overline{\alpha}\) attacks the assumption \(\alpha\) needed to apply the default rule, rather than deriving a competing conclusion. Assumption Introduction is therefore the transformation that allows ABA Learning to introduce defeasibility into the learnt framework.

Fact Subsumption removes redundant example-specific rules. If a rote-learnt rule is no longer needed because a more general rule already supports the relevant positive examples without causing negative examples to be accepted, the rote-learnt rule can be deleted. This prevents the learnt framework from retaining unnecessary facts once more general explanations have been found.

The distinction between rote and intensional learning is important. A purely rote solution may satisfy the examples by adding rules tied to individual constants. An intensional solution instead expresses reusable patterns through variables, background predicates, assumptions and contraries. In the context of this report, this distinction matters because the goal is to reason about learnable ABA structures, rather than only to memorise individual claims. The choice and order of transformation-rule applications can affect the learnt framework, which is why controlling the transformation process is an important part of applying ABA Learning in structured domains \cite{de2025greedy}.
